\section{Conclusions}
\label{sec:conclusions}

Direct edaphic cultivation in raw lunar regolith presents fundamental biological, toxicological, and fluid-mechanical barriers that cannot be resolved through chemical fertilization alone\cite{paul2022plants,ming1989lunar,lytvynenko2021lunar}. Molecular, transcriptomic, and toxicological investigations confirm that pristine lunar dust contains reactive nanophase metallic iron ($\mathrm{npFe}^0$), unpassivated silicate dangling bonds, and razor-sharp mineral shards that trigger chronic oxidative stress cascades, inflict 8-oxoguanine genomic lesions, accelerate telomere shortening, and arrest plant root development\cite{paul2022plants,ming1989lunar,hendrix2024reactivity,barbero2024telomere}. Simultaneously, the reduction of gravitational acceleration to $1/6\,g$ lowers the Bond number, establishing a capillary-dominated hydraulic regime in which polydisperse regolith beds experience severe capillary waterlogging and root-zone hypoxia\cite{ming1989lunar,ostrach1982low,podolskiy2005capillary}.

Transforming in-situ regolith into engineered sintered ceramics offers an ISRU manufacturing paradigm that fundamentally circumvents these edaphic constraints\cite{ming1989lunar,taylor2005microwave}. Thermal processing above $1050^\circ\mathrm{C}$ passivates reactive geochemical interfaces by solid-solubilizing $\mathrm{npFe}^0$ into stable augite and olivine crystal frameworks, rounds abrasive grain edges, and quenches dangling surface radicals, creating an inert, biocompatible growth substrate\cite{ming1989lunar,schoonen2020hydroxyl,xi2026multiphysics}. Pre-sintering classification of regolith particles to $45\text{--}75\,\mu\mathrm{m}$ produces monodisperse pore throat diameters of $30\text{--}60\,\mu\mathrm{m}$, stabilizing matric potential within the optimal window ($-2.5$ to $-5.0\,\mathrm{kPa}$) while maintaining air-filled porosity above $25\text{--}30\%$ for unhindered aerobic root respiration under fractional gravity\cite{ming1989lunar,heinse2007porous}.

From a mission systems perspective, the upfront launch mass required for an autonomous sintering plant is offset rapidly by operational efficiencies\cite{ming1989lunar,hanford2004advanced,metzger2023cost}. Life cycle assessments based on NASA Equivalent System Mass formulations demonstrate an operational breakeven horizon between $18$ and $22$ months compared to direct regolith cultivation (and as low as $9\text{--}15$ months under surface nuclear power architectures)\cite{ming1989lunar,chinnannan2025spirulina,colozza2020power}. This advantage is driven by $>98.5\%$ water recovery efficiency, intrinsic $\mathrm{pH}$ stability, and the elimination of chemical buffer resupply\cite{ming1989lunar,anderson2021lunar}. When linked with saprotrophic fungal mycoponics for inedible biomass upcycling, external microbial bioweathering reactors for closed-loop mineral replenishment, and cyclic $500^\circ\mathrm{C}$ pyrolytic bake-outs for substrate regeneration, sintered regolith ceramics establish a robust, durable foundation for sustained biogenerative human life support on the lunar surface\cite{ming1989lunar,liu2020bioregenerative,cockell2018biorock,barta1994thermal}.
